\documentclass[11pt,letterpaper]{article}

\usepackage[T1]{fontenc}
\usepackage[utf8]{inputenc}
\usepackage{lmodern}
\usepackage[margin=1in]{geometry}
\usepackage{microtype}
\usepackage{xcolor}
\usepackage{booktabs}
\usepackage{tabularx}
\usepackage{enumitem}
\usepackage{titlesec}
\usepackage{parskip}
\usepackage[hidelinks]{hyperref}

\definecolor{accent}{HTML}{3B6D11}
\definecolor{titledark}{HTML}{2C4A20}
\definecolor{muted}{HTML}{5F5E5A}
\definecolor{rulegray}{HTML}{C9C7BE}
\definecolor{headfill}{HTML}{E8EDE4}

\titleformat{\section}{\normalfont\large\bfseries\color{accent}}{}{0pt}{}
\titlespacing*{\section}{0pt}{1.4em}{0.5em}

\newcommand{\unit}[1]{\ensuremath{\mathrm{#1}}}
\newcommand{\Vcmax}{\ensuremath{V_{\mathrm{cmax}}}}
\newcommand{\Jmax}{\ensuremath{J_{\mathrm{max}}}}

\setlist[itemize]{leftmargin=1.4em, itemsep=0.35em, topsep=0.4em}
\newcolumntype{L}{>{\raggedright\arraybackslash}X}

\begin{document}

{\noindent\color{titledark}\LARGE\bfseries Spectrophotometer vs.\ spectroradiometer\par}
\vspace{0.4em}
{\noindent\color{rulegray}\rule{\textwidth}{0.8pt}\par}
\vspace{0.5em}
{\noindent\color{muted}\itshape Two instruments that both produce a spectrum, separated by where the
light comes from and what the numbers mean.\par}
\vspace{1.2em}

\noindent\textbf{In short:} a spectrophotometer makes its own light and measures how a sample
interacts with it. A spectroradiometer measures whatever light arrives, in absolute physical units.

\section{Spectrophotometer}

\begin{itemize}
  \item \textbf{Has an internal light source} --- tungsten-halogen, deuterium, or xenon flash ---
        plus a sample holder in a fixed, controlled geometry.
  \item \textbf{Measures a ratio, not absolute energy}: transmittance or reflectance relative to a
        reference (a blank cuvette, a white Spectralon standard) measured through the same optical
        path. Output is unitless, reported as $0$--$1$ or a percentage, or as absorbance
        $A = -\log_{10} T$.
  \item Because source, path length, and geometry are fixed, the ratio is stable and reproducible.
        That is the entire point of the design.
  \item \textbf{Typical instruments:} Cary 5000, Shimadzu UV-3600, and the integrating-sphere leaf
        clips such as the ASD or PSR contact probe and the LI-COR 6800 leaf spectrometer. Anything
        using an integrating sphere for leaf reflectance and transmittance is effectively operating
        in spectrophotometer mode.
\end{itemize}

\section{Spectroradiometer}

\begin{itemize}
  \item \textbf{No internal source} --- it is a passive detector pointed at a scene, calibrated to
        report absolute radiometric quantities: spectral radiance in
        $\unit{W\,m^{-2}\,sr^{-1}\,nm^{-1}}$, or irradiance in $\unit{W\,m^{-2}\,nm^{-1}}$.
  \item Requires a radiometric calibration traceable to a standard lamp, together with a defined
        field of view or a cosine-corrected diffuser.
  \item To obtain reflectance you take two measurements --- target radiance and a white reference
        panel under the same illumination --- and form the ratio yourself. Illumination is the sun,
        or whatever is there, so it changes between scans.
  \item \textbf{Typical instruments:} ASD FieldSpec, Spectra Vista SVC HR-1024i, PSR+, and the
        sensors aboard airborne and satellite imaging spectrometers such as AVIRIS-NG, EMIT, and
        PRISMA.
\end{itemize}

\section{Side by side}

\vspace{0.3em}
\noindent
{\small\renewcommand{\arraystretch}{1.15}
\begin{tabularx}{\textwidth}{@{}>{\raggedright\arraybackslash}p{0.21\textwidth} L L@{}}
\toprule
& \textbf{Spectrophotometer} & \textbf{Spectroradiometer} \\
\midrule
\textbf{Illumination}
  & Internal, constant, controlled
  & Ambient --- sun, sky, or whatever is present \\[0.4em]
\textbf{Native output}
  & A ratio: reflectance, transmittance, absorbance
  & Absolute radiance or irradiance in physical units \\[0.4em]
\textbf{Getting reflectance}
  & Direct, in one step against a blank or white standard
  & Two steps: target radiance ratioed to a white panel \\[0.4em]
\textbf{Scale of measurement}
  & Leaf or sample
  & Leaf through canopy to whole scene \\[0.4em]
\textbf{Dominant noise sources}
  & Instrument drift, stray light, sphere substitution error
  & Sun angle, cloud, BRDF, atmosphere, background \\
\bottomrule
\end{tabularx}
}
\vspace{0.6em}

\section{Why it matters for hyperspectral trait work}

For partial least squares regression models predicting nitrogen, LMA, chlorophyll, carotenoids,
\Vcmax, \Jmax, cellulose, and lignin, leaf-level spectra from a contact probe or integrating sphere
--- spectrophotometer-style, with a fixed artificial source --- are far cleaner than anything
measured passively. That is the main reason leaf-level trait models validate so much better than
canopy-level ones.

When the same model is transferred to canopy spectroradiometer data, the degradation you observe is
mostly illumination geometry, structural scattering, and background contribution. It is rarely a
failure of the underlying chemistry--spectra relationship.

\section{A note on terminology}

Field instruments such as the ASD FieldSpec are spectroradiometers. Attach a contact probe with its
own lamp, however, and you are using one as a spectrophotometer. The hardware category and the
measurement mode are not the same thing, and papers are not always careful about the distinction.

\end{document}
